% Generated draft: requires mathematical and source review.
\documentclass[11pt]{article}
\usepackage[margin=25mm]{geometry}
\usepackage{fontspec}
\setmainfont{Latin Modern Roman}
\usepackage{amsmath,amssymb,mathtools,unicode-math}
\setmathfont{Latin Modern Math}
\usepackage{enumitem,needspace,xurl,hyperref}
\setlength{\parindent}{0pt}
\setlength{\parskip}{5pt}
\setlength{\emergencystretch}{3em}
\newcommand{\N}{\mathbb{N}}
\newcommand{\Z}{\mathbb{Z}}
\newcommand{\Q}{\mathbb{Q}}
\newcommand{\R}{\mathbb{R}}
\newcommand{\C}{\mathbb{C}}
\newcommand{\F}{\mathbb{F}}
\newcommand{\A}{\mathbb{A}}
\newcommand{\PP}{\mathbb{P}}
\newcommand{\E}{\mathbb{E}}
\newcommand{\eps}{\varepsilon}
\newcommand{\dd}{\,\mathrm{d}}
\begin{document}
\section*{UnsolvedMath ID: 30004714}
\textbf{Source problem code:} catalogue entry
\subsection{Definitions and mathematical statement}
Let \(k\) be a field of characteristic \(p>0\), and let \(V\) be a finite-dimensional \(k\)-vector space of dimension \(n\geq 1\). Write \(G=\operatorname{GL}_k(V)\) and \(X=\operatorname{Flag}(V)\) for the complete flag variety. Its points over a field extension \(K/k\) are chains
\[
0=F_0\subset F_1\subset\cdots\subset F_n=V\otimes_k K,
\qquad \dim_K F_j=j.
\]
Let \(\mathcal S_j\) be the tautological rank-\(j\) subbundle of \(V\otimes_k\mathcal O_X\), and set \(\mathcal Q_j=\mathcal S_j/\mathcal S_{j-1}\). For \(\mathbf a=(a_1,\ldots,a_n)\in\Z^n\), define
\[
\mathcal L_{\mathbf a}=\bigotimes_{j=1}^{n}\mathcal Q_j^{\otimes a_j},
\]
where a negative tensor power means the corresponding positive tensor power of the dual. Equip these line bundles with the \(G\)-linearizations induced by the tautological action. Thus each sheaf cohomology group \(H^q(X,\mathcal L_{\mathbf a})\) is a finite-dimensional rational representation of the algebraic group \(G\) over \(k\).

With this explicit equivariant convention, the requested general problem is to determine, for every positive characteristic field \(k\), every such \(V\), every \(\mathbf a\in\Z^n\), and every integer \(q\geq0\), the isomorphism class of
\[
H^q\!\left(\operatorname{Flag}(V),\mathcal L_{\mathbf a}\right)
\]
as a rational \(\operatorname{GL}_k(V)\)-representation. The weaker requested target is a necessary and sufficient criterion, in these parameters, for
\[
H^q\!\left(\operatorname{Flag}(V),\mathcal L_{\mathbf a}\right)=0.
\]
Determining a representation here means determining its module isomorphism class; a character formula alone need not determine that class in positive characteristic.

The source does not specify a linearization or a weight-sign convention. The tautological construction above makes these choices explicit for this draft. An underlying line bundle alone does not specify its \(G\)-representation on cohomology: changing a linearization by a character of \(G\) changes the resulting action. In particular, adding the same integer to all entries of \(\mathbf a\) tensors the equivariant line bundle by a power of the determinant character, while leaving its underlying line-bundle isomorphism class unchanged. The source also does not explicitly specify the dimension range or whether algebraic closedness was intended; the displayed formulation retains arbitrary positive characteristic fields and makes the conventional finite-dimensional setting explicit.

The supplied dated curation labels the problem partially solved. It attributes vanishing and nonvanishing results and character formulas for restricted families to Gao--Raicu, and recursions and characteristic-two formulas to later work, while describing the full representation problem as open. The applicable families and hypotheses are not supplied, so this draft does not interpret those claims as resolving the general vanishing target above.
\subsection{Short English statement}
Determine the cohomology of every line bundle on a complete flag variety in positive characteristic, including its full \(\operatorname{GL}(V)\)-module structure, or at least give an exact criterion for vanishing. The draft uses explicitly defined tautological linearizations because the source leaves the equivariant convention unspecified. The supplied partially solved label and descriptions of prior results require independent mathematical review.
\subsection{Sources}
\begin{itemize}
\item Local registry: \texttt{data/unsolvedmath/problems.json}, permanent UnsolvedMath ID \(30004714\), source problem code \texttt{catalogue entry}, title \emph{Flag-Variety Line-Bundle Cohomology in Positive Characteristic}. The supplied record belongs to the Oberwolfach Reports Open Problems collection, records proposed year 2021, and names no individual proposer.
\item Supplied report citation: \emph{Classical Algebraic Geometry (2022)}, Oberwolfach Reports, DOI 10.4171/owr/2021/29, pp.\ 1519--1577. \url{https://doi.org/10.4171/owr/2021/29}. The supplied background describes these reports as documenting mathematical workshops and their open-problem discussions.
\item Supplied literature reference: Gao--Raicu, \url{https://arxiv.org/abs/2210.03656}. No paper title or precise theorem reference is supplied.
\item Supplied literature reference: Kyomuhangi--Raicu--Reed, \url{https://arxiv.org/abs/2411.13450}. No paper title or precise theorem reference is supplied.
\item Literature assessment dated exactly \texttt{2026-08-22}: the supplied curation records partial progress on vanishing, nonvanishing, restricted character formulas, recursions, and characteristic-two formulas, and describes the full representation problem as open. This is supplied dated curation derived from the OWR material, not an independently verified literature assessment.
\item The supplied statement status is \texttt{corrected\_verified}; its accompanying explanation says that a publication-format audit preserved mathematical content while repairing notation or typography. This supplied label does not establish independent mathematical verification of this draft.
\end{itemize}
\end{document}
